% =======================================================================
%   Shared style for CS 300 notes.
%   Usage: place `% =======================================================================
%   Shared style for CS 300 notes.
%   Usage: place `% =======================================================================
%   Shared style for CS 300 notes.
%   Usage: place `% =======================================================================
%   Shared style for CS 300 notes.
%   Usage: place `\input{../../notes-style.tex}` right after \documentclass
%   in each chapter's lectures.tex and notes.tex.
% =======================================================================

% ---- Layout -----------------------------------------------------------
\usepackage[margin=1in]{geometry}
\usepackage{parskip}        % space between paragraphs, no indent
\usepackage{microtype}      % subtle kerning/spacing improvements

% ---- Colors (muted palette, easy on light-sensitive eyes) -------------
\usepackage{xcolor}
\definecolor{pagebg}       {HTML}{FAF6F0}  % warm cream background
\definecolor{bodyfg}       {HTML}{3D3A36}  % warm dark gray body text
\definecolor{heading}      {HTML}{3D6A6B}  % muted teal
\definecolor{subheading}   {HTML}{5B7A9C}  % dusty blue
\definecolor{accent}       {HTML}{8A6B4A}  % warm tan
\definecolor{linkcol}      {HTML}{5B7A9C}  % same as subheading
\definecolor{mutedtext}    {HTML}{6B6560}  % secondary / caption text
\definecolor{rulecol}      {HTML}{C8BFB0}  % separator / rule

% Callout box palette
\definecolor{defnFrame}    {HTML}{9DB89B}  \definecolor{defnBack}    {HTML}{EEF2E8}
\definecolor{ideaFrame}    {HTML}{9FB4CD}  \definecolor{ideaBack}    {HTML}{E8EEF5}
\definecolor{warnFrame}    {HTML}{CFAE87}  \definecolor{warnBack}    {HTML}{F5EDDF}
\definecolor{noteFrame}    {HTML}{BFA6AE}  \definecolor{noteBack}    {HTML}{F2E8EE}
\definecolor{exFrame}      {HTML}{9FBDBD}  \definecolor{exBack}      {HTML}{E8F0F0}
\definecolor{codeBack}     {HTML}{F3EDE2}
\definecolor{codeBorder}   {HTML}{D4CBBA}
\definecolor{codeKeyword}  {HTML}{6B4F7A}
\definecolor{codeString}   {HTML}{8A6B4A}
\definecolor{codeComment}  {HTML}{8A857F}

% ---- Page background + base text color --------------------------------
\usepackage{pagecolor}
\pagecolor{pagebg}
\color{bodyfg}

% ---- Fonts ------------------------------------------------------------
\usepackage[T1]{fontenc}
\IfFileExists{libertine.sty}{%
    \usepackage{libertine}      % warm, readable serif
    \usepackage[scaled=0.95]{inconsolata}  % softer monospace
}{}

% ---- Hyperref ---------------------------------------------------------
\usepackage{hyperref}
\hypersetup{
    colorlinks=true,
    linkcolor=linkcol,
    urlcolor=linkcol,
    citecolor=linkcol,
    pdfborder={0 0 0},
}

% ---- Section styling --------------------------------------------------
\usepackage[explicit]{titlesec}
\titleformat{\section}
    {\color{heading}\Large\bfseries}
    {\color{heading}\thesection\hspace{0.6em}}{0pt}{#1}
\titleformat{\subsection}
    {\color{subheading}\large\bfseries}
    {\color{subheading}\thesubsection\hspace{0.5em}}{0pt}{#1}
\titleformat{\subsubsection}
    {\color{accent}\normalsize\bfseries}
    {\color{accent}\thesubsubsection\hspace{0.5em}}{0pt}{#1}
\titleformat{\paragraph}[runin]
    {\color{accent}\bfseries}{}{0pt}{#1.}

\titlespacing*{\section}{0pt}{1.6em}{0.6em}
\titlespacing*{\subsection}{0pt}{1.2em}{0.4em}

% ---- Enumerate / itemize spacing --------------------------------------
\usepackage{enumitem}
\setlist[itemize]{leftmargin=1.4em, itemsep=2pt, topsep=4pt}
\setlist[enumerate]{leftmargin=1.6em, itemsep=2pt, topsep=4pt}

% ---- Code listings ----------------------------------------------------
\usepackage{listings}
\lstdefinestyle{notescode}{
    backgroundcolor=\color{codeBack},
    basicstyle=\ttfamily\small\color{bodyfg},
    keywordstyle=\color{codeKeyword}\bfseries,
    stringstyle=\color{codeString},
    commentstyle=\color{codeComment}\itshape,
    frame=single,
    framesep=6pt,
    rulecolor=\color{codeBorder},
    showstringspaces=false,
    breaklines=true,
    xleftmargin=10pt,
    xrightmargin=10pt,
    aboveskip=8pt,
    belowskip=8pt,
}
\lstset{style=notescode}

% ---- Callout boxes (tcolorbox) ----------------------------------------
\usepackage[most]{tcolorbox}

\newtcolorbox{defnbox}[1][]{
    enhanced, breakable,
    colback=defnBack, colframe=defnFrame, coltitle=bodyfg,
    title={\bfseries Definition \textemdash\ #1},
    fonttitle=\bfseries,
    boxrule=0.5pt, arc=2pt, left=8pt, right=8pt, top=6pt, bottom=6pt,
    attach boxed title to top left={xshift=10pt, yshift=-6pt},
    boxed title style={colback=defnFrame, colframe=defnFrame, arc=2pt},
    before skip=8pt, after skip=8pt,
}
\newtcolorbox{ideabox}[1][]{
    enhanced, breakable,
    colback=ideaBack, colframe=ideaFrame, coltitle=bodyfg,
    title={\bfseries Key idea #1},
    boxrule=0.5pt, arc=2pt, left=8pt, right=8pt, top=6pt, bottom=6pt,
    attach boxed title to top left={xshift=10pt, yshift=-6pt},
    boxed title style={colback=ideaFrame, colframe=ideaFrame, arc=2pt},
    before skip=8pt, after skip=8pt,
}
\newtcolorbox{warnbox}[1][]{
    enhanced, breakable,
    colback=warnBack, colframe=warnFrame, coltitle=bodyfg,
    title={\bfseries Gotcha #1},
    boxrule=0.5pt, arc=2pt, left=8pt, right=8pt, top=6pt, bottom=6pt,
    attach boxed title to top left={xshift=10pt, yshift=-6pt},
    boxed title style={colback=warnFrame, colframe=warnFrame, arc=2pt},
    before skip=8pt, after skip=8pt,
}
\newtcolorbox{notebox}[1][]{
    enhanced, breakable,
    colback=noteBack, colframe=noteFrame, coltitle=bodyfg,
    title={\bfseries Aside #1},
    boxrule=0.5pt, arc=2pt, left=8pt, right=8pt, top=6pt, bottom=6pt,
    attach boxed title to top left={xshift=10pt, yshift=-6pt},
    boxed title style={colback=noteFrame, colframe=noteFrame, arc=2pt},
    before skip=8pt, after skip=8pt,
}
\newtcolorbox{examplebox}[1][]{
    enhanced, breakable,
    colback=exBack, colframe=exFrame, coltitle=bodyfg,
    title={\bfseries Example #1},
    boxrule=0.5pt, arc=2pt, left=8pt, right=8pt, top=6pt, bottom=6pt,
    attach boxed title to top left={xshift=10pt, yshift=-6pt},
    boxed title style={colback=exFrame, colframe=exFrame, arc=2pt},
    before skip=8pt, after skip=8pt,
}

% ---- TikZ graph styles (added 2026-04-25 for ch_10 augmentation) ------
% tcolorbox[most] already loads tikz; this block declares the libraries
% + shared `vertex` / `edge` styles used by ch_10's general directed-graph
% diagrams. ch_9's tree-style TikZ uses inline overrides and is
% unaffected. Simple circle nodes + arrow edges only -- no production
% styling.
\usetikzlibrary{arrows.meta, positioning, calc}
\tikzset{
    vertex/.style={circle, draw, minimum size=7mm, inner sep=0pt,
                   font=\small, thick},
    vsource/.style={vertex, fill=ideaBack},
    vdone/.style={vertex, fill=defnBack},
    edge/.style={->, >=Stealth, thick},
    uedge/.style={-, thick},
    treeedge/.style={->, >=Stealth, thick},
    backedge/.style={->, >=Stealth, thick, dashed, red!70!black},
    fwdedge/.style={->, >=Stealth, thick, blue!60!black},
    crossedge/.style={->, >=Stealth, thick, green!50!black, dotted},
    mstedge/.style={-, very thick, accent},
}

% ---- Title block ------------------------------------------------------
\usepackage{titling}
\pretitle{\begin{center}\color{heading}\LARGE\bfseries}
\posttitle{\end{center}\vskip -0.4em}
\preauthor{\begin{center}\color{mutedtext}\normalsize}
\postauthor{\end{center}}
\predate{\begin{center}\color{mutedtext}\small}
\postdate{\end{center}\vspace{0.4em}{\color{rulecol}\hrule height 0.4pt}\vspace{1.2em}}
` right after \documentclass
%   in each chapter's lectures.tex and notes.tex.
% =======================================================================

% ---- Layout -----------------------------------------------------------
\usepackage[margin=1in]{geometry}
\usepackage{parskip}        % space between paragraphs, no indent
\usepackage{microtype}      % subtle kerning/spacing improvements

% ---- Colors (muted palette, easy on light-sensitive eyes) -------------
\usepackage{xcolor}
\definecolor{pagebg}       {HTML}{FAF6F0}  % warm cream background
\definecolor{bodyfg}       {HTML}{3D3A36}  % warm dark gray body text
\definecolor{heading}      {HTML}{3D6A6B}  % muted teal
\definecolor{subheading}   {HTML}{5B7A9C}  % dusty blue
\definecolor{accent}       {HTML}{8A6B4A}  % warm tan
\definecolor{linkcol}      {HTML}{5B7A9C}  % same as subheading
\definecolor{mutedtext}    {HTML}{6B6560}  % secondary / caption text
\definecolor{rulecol}      {HTML}{C8BFB0}  % separator / rule

% Callout box palette
\definecolor{defnFrame}    {HTML}{9DB89B}  \definecolor{defnBack}    {HTML}{EEF2E8}
\definecolor{ideaFrame}    {HTML}{9FB4CD}  \definecolor{ideaBack}    {HTML}{E8EEF5}
\definecolor{warnFrame}    {HTML}{CFAE87}  \definecolor{warnBack}    {HTML}{F5EDDF}
\definecolor{noteFrame}    {HTML}{BFA6AE}  \definecolor{noteBack}    {HTML}{F2E8EE}
\definecolor{exFrame}      {HTML}{9FBDBD}  \definecolor{exBack}      {HTML}{E8F0F0}
\definecolor{codeBack}     {HTML}{F3EDE2}
\definecolor{codeBorder}   {HTML}{D4CBBA}
\definecolor{codeKeyword}  {HTML}{6B4F7A}
\definecolor{codeString}   {HTML}{8A6B4A}
\definecolor{codeComment}  {HTML}{8A857F}

% ---- Page background + base text color --------------------------------
\usepackage{pagecolor}
\pagecolor{pagebg}
\color{bodyfg}

% ---- Fonts ------------------------------------------------------------
\usepackage[T1]{fontenc}
\IfFileExists{libertine.sty}{%
    \usepackage{libertine}      % warm, readable serif
    \usepackage[scaled=0.95]{inconsolata}  % softer monospace
}{}

% ---- Hyperref ---------------------------------------------------------
\usepackage{hyperref}
\hypersetup{
    colorlinks=true,
    linkcolor=linkcol,
    urlcolor=linkcol,
    citecolor=linkcol,
    pdfborder={0 0 0},
}

% ---- Section styling --------------------------------------------------
\usepackage[explicit]{titlesec}
\titleformat{\section}
    {\color{heading}\Large\bfseries}
    {\color{heading}\thesection\hspace{0.6em}}{0pt}{#1}
\titleformat{\subsection}
    {\color{subheading}\large\bfseries}
    {\color{subheading}\thesubsection\hspace{0.5em}}{0pt}{#1}
\titleformat{\subsubsection}
    {\color{accent}\normalsize\bfseries}
    {\color{accent}\thesubsubsection\hspace{0.5em}}{0pt}{#1}
\titleformat{\paragraph}[runin]
    {\color{accent}\bfseries}{}{0pt}{#1.}

\titlespacing*{\section}{0pt}{1.6em}{0.6em}
\titlespacing*{\subsection}{0pt}{1.2em}{0.4em}

% ---- Enumerate / itemize spacing --------------------------------------
\usepackage{enumitem}
\setlist[itemize]{leftmargin=1.4em, itemsep=2pt, topsep=4pt}
\setlist[enumerate]{leftmargin=1.6em, itemsep=2pt, topsep=4pt}

% ---- Code listings ----------------------------------------------------
\usepackage{listings}
\lstdefinestyle{notescode}{
    backgroundcolor=\color{codeBack},
    basicstyle=\ttfamily\small\color{bodyfg},
    keywordstyle=\color{codeKeyword}\bfseries,
    stringstyle=\color{codeString},
    commentstyle=\color{codeComment}\itshape,
    frame=single,
    framesep=6pt,
    rulecolor=\color{codeBorder},
    showstringspaces=false,
    breaklines=true,
    xleftmargin=10pt,
    xrightmargin=10pt,
    aboveskip=8pt,
    belowskip=8pt,
}
\lstset{style=notescode}

% ---- Callout boxes (tcolorbox) ----------------------------------------
\usepackage[most]{tcolorbox}

\newtcolorbox{defnbox}[1][]{
    enhanced, breakable,
    colback=defnBack, colframe=defnFrame, coltitle=bodyfg,
    title={\bfseries Definition \textemdash\ #1},
    fonttitle=\bfseries,
    boxrule=0.5pt, arc=2pt, left=8pt, right=8pt, top=6pt, bottom=6pt,
    attach boxed title to top left={xshift=10pt, yshift=-6pt},
    boxed title style={colback=defnFrame, colframe=defnFrame, arc=2pt},
    before skip=8pt, after skip=8pt,
}
\newtcolorbox{ideabox}[1][]{
    enhanced, breakable,
    colback=ideaBack, colframe=ideaFrame, coltitle=bodyfg,
    title={\bfseries Key idea #1},
    boxrule=0.5pt, arc=2pt, left=8pt, right=8pt, top=6pt, bottom=6pt,
    attach boxed title to top left={xshift=10pt, yshift=-6pt},
    boxed title style={colback=ideaFrame, colframe=ideaFrame, arc=2pt},
    before skip=8pt, after skip=8pt,
}
\newtcolorbox{warnbox}[1][]{
    enhanced, breakable,
    colback=warnBack, colframe=warnFrame, coltitle=bodyfg,
    title={\bfseries Gotcha #1},
    boxrule=0.5pt, arc=2pt, left=8pt, right=8pt, top=6pt, bottom=6pt,
    attach boxed title to top left={xshift=10pt, yshift=-6pt},
    boxed title style={colback=warnFrame, colframe=warnFrame, arc=2pt},
    before skip=8pt, after skip=8pt,
}
\newtcolorbox{notebox}[1][]{
    enhanced, breakable,
    colback=noteBack, colframe=noteFrame, coltitle=bodyfg,
    title={\bfseries Aside #1},
    boxrule=0.5pt, arc=2pt, left=8pt, right=8pt, top=6pt, bottom=6pt,
    attach boxed title to top left={xshift=10pt, yshift=-6pt},
    boxed title style={colback=noteFrame, colframe=noteFrame, arc=2pt},
    before skip=8pt, after skip=8pt,
}
\newtcolorbox{examplebox}[1][]{
    enhanced, breakable,
    colback=exBack, colframe=exFrame, coltitle=bodyfg,
    title={\bfseries Example #1},
    boxrule=0.5pt, arc=2pt, left=8pt, right=8pt, top=6pt, bottom=6pt,
    attach boxed title to top left={xshift=10pt, yshift=-6pt},
    boxed title style={colback=exFrame, colframe=exFrame, arc=2pt},
    before skip=8pt, after skip=8pt,
}

% ---- TikZ graph styles (added 2026-04-25 for ch_10 augmentation) ------
% tcolorbox[most] already loads tikz; this block declares the libraries
% + shared `vertex` / `edge` styles used by ch_10's general directed-graph
% diagrams. ch_9's tree-style TikZ uses inline overrides and is
% unaffected. Simple circle nodes + arrow edges only -- no production
% styling.
\usetikzlibrary{arrows.meta, positioning, calc}
\tikzset{
    vertex/.style={circle, draw, minimum size=7mm, inner sep=0pt,
                   font=\small, thick},
    vsource/.style={vertex, fill=ideaBack},
    vdone/.style={vertex, fill=defnBack},
    edge/.style={->, >=Stealth, thick},
    uedge/.style={-, thick},
    treeedge/.style={->, >=Stealth, thick},
    backedge/.style={->, >=Stealth, thick, dashed, red!70!black},
    fwdedge/.style={->, >=Stealth, thick, blue!60!black},
    crossedge/.style={->, >=Stealth, thick, green!50!black, dotted},
    mstedge/.style={-, very thick, accent},
}

% ---- Title block ------------------------------------------------------
\usepackage{titling}
\pretitle{\begin{center}\color{heading}\LARGE\bfseries}
\posttitle{\end{center}\vskip -0.4em}
\preauthor{\begin{center}\color{mutedtext}\normalsize}
\postauthor{\end{center}}
\predate{\begin{center}\color{mutedtext}\small}
\postdate{\end{center}\vspace{0.4em}{\color{rulecol}\hrule height 0.4pt}\vspace{1.2em}}
` right after \documentclass
%   in each chapter's lectures.tex and notes.tex.
% =======================================================================

% ---- Layout -----------------------------------------------------------
\usepackage[margin=1in]{geometry}
\usepackage{parskip}        % space between paragraphs, no indent
\usepackage{microtype}      % subtle kerning/spacing improvements

% ---- Colors (muted palette, easy on light-sensitive eyes) -------------
\usepackage{xcolor}
\definecolor{pagebg}       {HTML}{FAF6F0}  % warm cream background
\definecolor{bodyfg}       {HTML}{3D3A36}  % warm dark gray body text
\definecolor{heading}      {HTML}{3D6A6B}  % muted teal
\definecolor{subheading}   {HTML}{5B7A9C}  % dusty blue
\definecolor{accent}       {HTML}{8A6B4A}  % warm tan
\definecolor{linkcol}      {HTML}{5B7A9C}  % same as subheading
\definecolor{mutedtext}    {HTML}{6B6560}  % secondary / caption text
\definecolor{rulecol}      {HTML}{C8BFB0}  % separator / rule

% Callout box palette
\definecolor{defnFrame}    {HTML}{9DB89B}  \definecolor{defnBack}    {HTML}{EEF2E8}
\definecolor{ideaFrame}    {HTML}{9FB4CD}  \definecolor{ideaBack}    {HTML}{E8EEF5}
\definecolor{warnFrame}    {HTML}{CFAE87}  \definecolor{warnBack}    {HTML}{F5EDDF}
\definecolor{noteFrame}    {HTML}{BFA6AE}  \definecolor{noteBack}    {HTML}{F2E8EE}
\definecolor{exFrame}      {HTML}{9FBDBD}  \definecolor{exBack}      {HTML}{E8F0F0}
\definecolor{codeBack}     {HTML}{F3EDE2}
\definecolor{codeBorder}   {HTML}{D4CBBA}
\definecolor{codeKeyword}  {HTML}{6B4F7A}
\definecolor{codeString}   {HTML}{8A6B4A}
\definecolor{codeComment}  {HTML}{8A857F}

% ---- Page background + base text color --------------------------------
\usepackage{pagecolor}
\pagecolor{pagebg}
\color{bodyfg}

% ---- Fonts ------------------------------------------------------------
\usepackage[T1]{fontenc}
\IfFileExists{libertine.sty}{%
    \usepackage{libertine}      % warm, readable serif
    \usepackage[scaled=0.95]{inconsolata}  % softer monospace
}{}

% ---- Hyperref ---------------------------------------------------------
\usepackage{hyperref}
\hypersetup{
    colorlinks=true,
    linkcolor=linkcol,
    urlcolor=linkcol,
    citecolor=linkcol,
    pdfborder={0 0 0},
}

% ---- Section styling --------------------------------------------------
\usepackage[explicit]{titlesec}
\titleformat{\section}
    {\color{heading}\Large\bfseries}
    {\color{heading}\thesection\hspace{0.6em}}{0pt}{#1}
\titleformat{\subsection}
    {\color{subheading}\large\bfseries}
    {\color{subheading}\thesubsection\hspace{0.5em}}{0pt}{#1}
\titleformat{\subsubsection}
    {\color{accent}\normalsize\bfseries}
    {\color{accent}\thesubsubsection\hspace{0.5em}}{0pt}{#1}
\titleformat{\paragraph}[runin]
    {\color{accent}\bfseries}{}{0pt}{#1.}

\titlespacing*{\section}{0pt}{1.6em}{0.6em}
\titlespacing*{\subsection}{0pt}{1.2em}{0.4em}

% ---- Enumerate / itemize spacing --------------------------------------
\usepackage{enumitem}
\setlist[itemize]{leftmargin=1.4em, itemsep=2pt, topsep=4pt}
\setlist[enumerate]{leftmargin=1.6em, itemsep=2pt, topsep=4pt}

% ---- Code listings ----------------------------------------------------
\usepackage{listings}
\lstdefinestyle{notescode}{
    backgroundcolor=\color{codeBack},
    basicstyle=\ttfamily\small\color{bodyfg},
    keywordstyle=\color{codeKeyword}\bfseries,
    stringstyle=\color{codeString},
    commentstyle=\color{codeComment}\itshape,
    frame=single,
    framesep=6pt,
    rulecolor=\color{codeBorder},
    showstringspaces=false,
    breaklines=true,
    xleftmargin=10pt,
    xrightmargin=10pt,
    aboveskip=8pt,
    belowskip=8pt,
}
\lstset{style=notescode}

% ---- Callout boxes (tcolorbox) ----------------------------------------
\usepackage[most]{tcolorbox}

\newtcolorbox{defnbox}[1][]{
    enhanced, breakable,
    colback=defnBack, colframe=defnFrame, coltitle=bodyfg,
    title={\bfseries Definition \textemdash\ #1},
    fonttitle=\bfseries,
    boxrule=0.5pt, arc=2pt, left=8pt, right=8pt, top=6pt, bottom=6pt,
    attach boxed title to top left={xshift=10pt, yshift=-6pt},
    boxed title style={colback=defnFrame, colframe=defnFrame, arc=2pt},
    before skip=8pt, after skip=8pt,
}
\newtcolorbox{ideabox}[1][]{
    enhanced, breakable,
    colback=ideaBack, colframe=ideaFrame, coltitle=bodyfg,
    title={\bfseries Key idea #1},
    boxrule=0.5pt, arc=2pt, left=8pt, right=8pt, top=6pt, bottom=6pt,
    attach boxed title to top left={xshift=10pt, yshift=-6pt},
    boxed title style={colback=ideaFrame, colframe=ideaFrame, arc=2pt},
    before skip=8pt, after skip=8pt,
}
\newtcolorbox{warnbox}[1][]{
    enhanced, breakable,
    colback=warnBack, colframe=warnFrame, coltitle=bodyfg,
    title={\bfseries Gotcha #1},
    boxrule=0.5pt, arc=2pt, left=8pt, right=8pt, top=6pt, bottom=6pt,
    attach boxed title to top left={xshift=10pt, yshift=-6pt},
    boxed title style={colback=warnFrame, colframe=warnFrame, arc=2pt},
    before skip=8pt, after skip=8pt,
}
\newtcolorbox{notebox}[1][]{
    enhanced, breakable,
    colback=noteBack, colframe=noteFrame, coltitle=bodyfg,
    title={\bfseries Aside #1},
    boxrule=0.5pt, arc=2pt, left=8pt, right=8pt, top=6pt, bottom=6pt,
    attach boxed title to top left={xshift=10pt, yshift=-6pt},
    boxed title style={colback=noteFrame, colframe=noteFrame, arc=2pt},
    before skip=8pt, after skip=8pt,
}
\newtcolorbox{examplebox}[1][]{
    enhanced, breakable,
    colback=exBack, colframe=exFrame, coltitle=bodyfg,
    title={\bfseries Example #1},
    boxrule=0.5pt, arc=2pt, left=8pt, right=8pt, top=6pt, bottom=6pt,
    attach boxed title to top left={xshift=10pt, yshift=-6pt},
    boxed title style={colback=exFrame, colframe=exFrame, arc=2pt},
    before skip=8pt, after skip=8pt,
}

% ---- TikZ graph styles (added 2026-04-25 for ch_10 augmentation) ------
% tcolorbox[most] already loads tikz; this block declares the libraries
% + shared `vertex` / `edge` styles used by ch_10's general directed-graph
% diagrams. ch_9's tree-style TikZ uses inline overrides and is
% unaffected. Simple circle nodes + arrow edges only -- no production
% styling.
\usetikzlibrary{arrows.meta, positioning, calc}
\tikzset{
    vertex/.style={circle, draw, minimum size=7mm, inner sep=0pt,
                   font=\small, thick},
    vsource/.style={vertex, fill=ideaBack},
    vdone/.style={vertex, fill=defnBack},
    edge/.style={->, >=Stealth, thick},
    uedge/.style={-, thick},
    treeedge/.style={->, >=Stealth, thick},
    backedge/.style={->, >=Stealth, thick, dashed, red!70!black},
    fwdedge/.style={->, >=Stealth, thick, blue!60!black},
    crossedge/.style={->, >=Stealth, thick, green!50!black, dotted},
    mstedge/.style={-, very thick, accent},
}

% ---- Title block ------------------------------------------------------
\usepackage{titling}
\pretitle{\begin{center}\color{heading}\LARGE\bfseries}
\posttitle{\end{center}\vskip -0.4em}
\preauthor{\begin{center}\color{mutedtext}\normalsize}
\postauthor{\end{center}}
\predate{\begin{center}\color{mutedtext}\small}
\postdate{\end{center}\vspace{0.4em}{\color{rulecol}\hrule height 0.4pt}\vspace{1.2em}}
` right after \documentclass
%   in each chapter's lectures.tex and notes.tex.
% =======================================================================

% ---- Layout -----------------------------------------------------------
\usepackage[margin=1in]{geometry}
\usepackage{parskip}        % space between paragraphs, no indent
\usepackage{microtype}      % subtle kerning/spacing improvements

% ---- Colors (muted palette, easy on light-sensitive eyes) -------------
\usepackage{xcolor}
\definecolor{pagebg}       {HTML}{FAF6F0}  % warm cream background
\definecolor{bodyfg}       {HTML}{3D3A36}  % warm dark gray body text
\definecolor{heading}      {HTML}{3D6A6B}  % muted teal
\definecolor{subheading}   {HTML}{5B7A9C}  % dusty blue
\definecolor{accent}       {HTML}{8A6B4A}  % warm tan
\definecolor{linkcol}      {HTML}{5B7A9C}  % same as subheading
\definecolor{mutedtext}    {HTML}{6B6560}  % secondary / caption text
\definecolor{rulecol}      {HTML}{C8BFB0}  % separator / rule

% Callout box palette
\definecolor{defnFrame}    {HTML}{9DB89B}  \definecolor{defnBack}    {HTML}{EEF2E8}
\definecolor{ideaFrame}    {HTML}{9FB4CD}  \definecolor{ideaBack}    {HTML}{E8EEF5}
\definecolor{warnFrame}    {HTML}{CFAE87}  \definecolor{warnBack}    {HTML}{F5EDDF}
\definecolor{noteFrame}    {HTML}{BFA6AE}  \definecolor{noteBack}    {HTML}{F2E8EE}
\definecolor{exFrame}      {HTML}{9FBDBD}  \definecolor{exBack}      {HTML}{E8F0F0}
\definecolor{codeBack}     {HTML}{F3EDE2}
\definecolor{codeBorder}   {HTML}{D4CBBA}
\definecolor{codeKeyword}  {HTML}{6B4F7A}
\definecolor{codeString}   {HTML}{8A6B4A}
\definecolor{codeComment}  {HTML}{8A857F}

% ---- Page background + base text color --------------------------------
\usepackage{pagecolor}
\pagecolor{pagebg}
\color{bodyfg}

% ---- Fonts ------------------------------------------------------------
\usepackage[T1]{fontenc}
\IfFileExists{libertine.sty}{%
    \usepackage{libertine}      % warm, readable serif
    \usepackage[scaled=0.95]{inconsolata}  % softer monospace
}{}

% ---- Hyperref ---------------------------------------------------------
\usepackage{hyperref}
\hypersetup{
    colorlinks=true,
    linkcolor=linkcol,
    urlcolor=linkcol,
    citecolor=linkcol,
    pdfborder={0 0 0},
}

% ---- Section styling --------------------------------------------------
\usepackage[explicit]{titlesec}
\titleformat{\section}
    {\color{heading}\Large\bfseries}
    {\color{heading}\thesection\hspace{0.6em}}{0pt}{#1}
\titleformat{\subsection}
    {\color{subheading}\large\bfseries}
    {\color{subheading}\thesubsection\hspace{0.5em}}{0pt}{#1}
\titleformat{\subsubsection}
    {\color{accent}\normalsize\bfseries}
    {\color{accent}\thesubsubsection\hspace{0.5em}}{0pt}{#1}
\titleformat{\paragraph}[runin]
    {\color{accent}\bfseries}{}{0pt}{#1.}

\titlespacing*{\section}{0pt}{1.6em}{0.6em}
\titlespacing*{\subsection}{0pt}{1.2em}{0.4em}

% ---- Enumerate / itemize spacing --------------------------------------
\usepackage{enumitem}
\setlist[itemize]{leftmargin=1.4em, itemsep=2pt, topsep=4pt}
\setlist[enumerate]{leftmargin=1.6em, itemsep=2pt, topsep=4pt}

% ---- Code listings ----------------------------------------------------
\usepackage{listings}
\lstdefinestyle{notescode}{
    backgroundcolor=\color{codeBack},
    basicstyle=\ttfamily\small\color{bodyfg},
    keywordstyle=\color{codeKeyword}\bfseries,
    stringstyle=\color{codeString},
    commentstyle=\color{codeComment}\itshape,
    frame=single,
    framesep=6pt,
    rulecolor=\color{codeBorder},
    showstringspaces=false,
    breaklines=true,
    xleftmargin=10pt,
    xrightmargin=10pt,
    aboveskip=8pt,
    belowskip=8pt,
}
\lstset{style=notescode}

% ---- Callout boxes (tcolorbox) ----------------------------------------
\usepackage[most]{tcolorbox}

\newtcolorbox{defnbox}[1][]{
    enhanced, breakable,
    colback=defnBack, colframe=defnFrame, coltitle=bodyfg,
    title={\bfseries Definition \textemdash\ #1},
    fonttitle=\bfseries,
    boxrule=0.5pt, arc=2pt, left=8pt, right=8pt, top=6pt, bottom=6pt,
    attach boxed title to top left={xshift=10pt, yshift=-6pt},
    boxed title style={colback=defnFrame, colframe=defnFrame, arc=2pt},
    before skip=8pt, after skip=8pt,
}
\newtcolorbox{ideabox}[1][]{
    enhanced, breakable,
    colback=ideaBack, colframe=ideaFrame, coltitle=bodyfg,
    title={\bfseries Key idea #1},
    boxrule=0.5pt, arc=2pt, left=8pt, right=8pt, top=6pt, bottom=6pt,
    attach boxed title to top left={xshift=10pt, yshift=-6pt},
    boxed title style={colback=ideaFrame, colframe=ideaFrame, arc=2pt},
    before skip=8pt, after skip=8pt,
}
\newtcolorbox{warnbox}[1][]{
    enhanced, breakable,
    colback=warnBack, colframe=warnFrame, coltitle=bodyfg,
    title={\bfseries Gotcha #1},
    boxrule=0.5pt, arc=2pt, left=8pt, right=8pt, top=6pt, bottom=6pt,
    attach boxed title to top left={xshift=10pt, yshift=-6pt},
    boxed title style={colback=warnFrame, colframe=warnFrame, arc=2pt},
    before skip=8pt, after skip=8pt,
}
\newtcolorbox{notebox}[1][]{
    enhanced, breakable,
    colback=noteBack, colframe=noteFrame, coltitle=bodyfg,
    title={\bfseries Aside #1},
    boxrule=0.5pt, arc=2pt, left=8pt, right=8pt, top=6pt, bottom=6pt,
    attach boxed title to top left={xshift=10pt, yshift=-6pt},
    boxed title style={colback=noteFrame, colframe=noteFrame, arc=2pt},
    before skip=8pt, after skip=8pt,
}
\newtcolorbox{examplebox}[1][]{
    enhanced, breakable,
    colback=exBack, colframe=exFrame, coltitle=bodyfg,
    title={\bfseries Example #1},
    boxrule=0.5pt, arc=2pt, left=8pt, right=8pt, top=6pt, bottom=6pt,
    attach boxed title to top left={xshift=10pt, yshift=-6pt},
    boxed title style={colback=exFrame, colframe=exFrame, arc=2pt},
    before skip=8pt, after skip=8pt,
}

% ---- TikZ graph styles (added 2026-04-25 for ch_10 augmentation) ------
% tcolorbox[most] already loads tikz; this block declares the libraries
% + shared `vertex` / `edge` styles used by ch_10's general directed-graph
% diagrams. ch_9's tree-style TikZ uses inline overrides and is
% unaffected. Simple circle nodes + arrow edges only -- no production
% styling.
\usetikzlibrary{arrows.meta, positioning, calc}
\tikzset{
    vertex/.style={circle, draw, minimum size=7mm, inner sep=0pt,
                   font=\small, thick},
    vsource/.style={vertex, fill=ideaBack},
    vdone/.style={vertex, fill=defnBack},
    edge/.style={->, >=Stealth, thick},
    uedge/.style={-, thick},
    treeedge/.style={->, >=Stealth, thick},
    backedge/.style={->, >=Stealth, thick, dashed, red!70!black},
    fwdedge/.style={->, >=Stealth, thick, blue!60!black},
    crossedge/.style={->, >=Stealth, thick, green!50!black, dotted},
    mstedge/.style={-, very thick, accent},
}

% ---- Title block ------------------------------------------------------
\usepackage{titling}
\pretitle{\begin{center}\color{heading}\LARGE\bfseries}
\posttitle{\end{center}\vskip -0.4em}
\preauthor{\begin{center}\color{mutedtext}\normalsize}
\postauthor{\end{center}}
\predate{\begin{center}\color{mutedtext}\small}
\postdate{\end{center}\vspace{0.4em}{\color{rulecol}\hrule height 0.4pt}\vspace{1.2em}}
